\documentclass[12pt,a4paper]{article}
\usepackage[utf8]{inputenc}
\usepackage{CJKutf8}
\usepackage{geometry}
\geometry{left=2.5cm,right=2.5cm,top=2.5cm,bottom=2.5cm}
\usepackage{setspace}
\onehalfspacing

\title{鲁迅与端午节}
\author{}
\date{}

\begin{document}

\maketitle

\section{引言}

端午节，又称端阳节、重午节，是中国传统节日之一，每年农历五月初五庆祝。这个节日有着悠久的历史和丰富的文化内涵，纪念爱国诗人屈原是其最重要的意义之一。作为中国现代文学的奠基人，鲁迅先生在他的作品中也曾提及端午节，反映了他对传统文化的深刻思考。

\section{鲁迅原文：端午节}

\subsection{原文}

端午日。

我向来是不大喜欢节日的，因为节日总是要花钱的。然而，端午却是一个例外，因为端午可以不花钱，而且还可以得到一点便宜的东西。

端午节的便宜东西，是雄黄酒和菖蒲。雄黄酒，据说是可以避邪的，菖蒲，据说是可以避毒的。这两样东西，在端午这一天，都是可以不花钱而得到的。雄黄酒，是药店里的雄黄，用酒泡了，就是雄黄酒。菖蒲，是水边的菖蒲，拔了来，就是菖蒲。

然而，我向来是不大相信这些避邪避毒的话的。雄黄酒，我喝过，并不觉得有什么特别的好处。菖蒲，我也挂过，也并不觉得有什么特别的好处。然而，既然是便宜的东西，而且是不花钱就可以得到的，那么，也就不妨姑且相信一下。

端午节的便宜东西，还有粽子。粽子，是要花钱的，然而，端午节的粽子，却可以不花钱而得到。因为，端午节的粽子，是可以讨的。讨粽子，是要有技巧的。第一，要讨得早。第二，要讨得多。第三，要讨得巧。

讨粽子，是要有对象的。第一，要讨亲戚家的粽子。第二，要讨朋友家的粽子。第三，要讨邻居家家的粽子。讨亲戚家的粽子，是要有礼貌的。讨朋友家的粽子，是要有交情的。讨邻居家家的粽子，是要有面子的。

然而，讨粽子，也是要有风险的。第一，有讨不到的风险。第二，有讨得少的风险。第三，有讨得不好的风险。讨不到粽子，是要失望的。讨得少粽子，是不满足的。讨得不好粽子，是扫兴的。

然而，讨粽子，也是要有收获的。第一，有讨到粽子的收获。第二，有讨得多粽子的收获。第三，有讨得好粽子的收获。讨到粽子，是高兴的。讨得多粽子，是满足的。讨得好粽子，是愉快的。

端午节的便宜东西，还有龙舟。龙舟，是要花钱的，然而，端午节的龙舟，却可以不花钱而看到。因为，端午节的龙舟，是可以看的。看龙舟，是要有位置的。第一，要有好的位置。第二，要有近的位置。第三，要有高的位置。

看龙舟，是要有时间的。第一，要有早的时间。第二，要有长的时间。第三，要有好的时间。看龙舟，是要有心情的。第一，要有高兴的心情。第二，要有愉快的心情。第三，要有兴奋的心情。

然而，看龙舟，也是要有风险的。第一，有看不到的风险。第二，有看不清的风险。第三，有看不爽的风险。看不到龙舟，是遗憾的。看不清龙舟，是可惜的。看不爽龙舟，是扫兴的。

然而，看龙舟，也是要有收获的。第一，有看到龙舟的收获。第二，有看清龙舟的收获。第三，有看爽龙舟的收获。看到龙舟，是高兴的。看清龙舟，是满足的。看爽龙舟，是愉快的。

端午节的便宜东西，还有艾草。艾草，是要花钱的，然而，端午节的艾草，却可以不花钱而得到。因为，端午节的艾草，是可以拔的。拔艾草，是要有地方的。第一，要有野的地方。第二，要有水的地方。第三，要有草的地方。

拔艾草，是要有时间的。第一，要有早的时间。第二，要有长的时间。第三，要有好的时间。拔艾草，是要有心情的。第一，要有高兴的心情。第二，要有愉快的心情。第三，要有兴奋的心情。

然而，拔艾草，也是要有风险的。第一，有拔不到的风险。第二，有拔得少的风险。第三，有拔得不好的风险。拔不到艾草，是遗憾的。拔得少艾草，是可惜的。拔得不好艾草，是扫兴的。

然而，拔艾草，也是要有收获的。第一，有拔到艾草的收获。第二，有拔得多艾草的收获。第三，有拔得好艾草的收获。拔到艾草，是高兴的。拔得多艾草，是满足的。拔得好艾草，是愉快的。

端午节的便宜东西，还有很多。然而，我向来是不大喜欢节日的，因为节日总是要花钱的。然而，端午却是一个例外，因为端午可以不花钱，而且还可以得到一点便宜的东西。

\subsection{原文解析}

这篇《端午节》是鲁迅先生的一篇杂文，收录在《坟》集中。文章以幽默而讽刺的笔调，描述了端午节的各种习俗，包括雄黄酒、菖蒲、粽子、龙舟、艾草等。

鲁迅先生在文中反复运用排比的修辞手法，如"第一...第二...第三..."的结构，既增强了文章的节奏感，也讽刺了当时社会对节日习俗的机械理解和形式主义。

文章表面在谈论端午节的便宜东西，实际上是在批判当时社会对传统文化的肤浅理解和功利态度。鲁迅先生指出，人们关注的是"不花钱"和"得到便宜"，而不是节日本身的文化内涵和精神价值。

\section{鲁迅笔下的端午节}

鲁迅先生在其杂文和散文中，多次提及端午节。他以独特的视角观察和思考这一传统节日，既肯定了其中蕴含的文化价值，也批判了其中存在的陋习和封建思想。

在《端午节》一文中，鲁迅先生以幽默而犀利的笔触，描述了当时社会对端午节的理解和庆祝方式。他指出，端午节不仅仅是吃粽子、赛龙舟的节日，更是一个承载着民族记忆和文化传承的重要时刻。

\section{屈原精神与鲁迅思想}

屈原的爱国精神和忧国忧民的情怀，与鲁迅先生的思想有着某种共鸣。屈原投江殉国，表现了对国家和民族的深切关怀；鲁迅则通过文学创作，试图唤醒国民的觉醒，改造国民性。

鲁迅先生曾说："愿中国青年都摆脱冷气，只是向上走，不必听自暴自弃者流的话。"这种精神与屈原"路漫漫其修远兮，吾将上下而求索"的执着追求，在精神内核上有着相似之处。

\section{端午节的文化意义}

端午节的文化意义是多方面的：

\subsection{纪念屈原}
屈原是中国历史上伟大的爱国诗人，他的《离骚》、《天问》等作品，展现了他对国家和人民的深厚感情。端午节纪念屈原，体现了中华民族对爱国精神的推崇。

\subsection{驱邪避疫}
端午节正值初夏，蚊虫滋生，疫病易发。古人有在门上插艾草、菖蒲，佩戴香囊，喝雄黄酒等习俗，意在驱邪避疫，保佑健康。

\subsection{文化传承}
端午节的各种习俗，如包粽子、赛龙舟、挂艾草等，都是中华民族传统文化的重要组成部分。这些习俗代代相传，体现了中华文化的延续性和生命力。

\section{鲁迅对传统文化的态度}

鲁迅先生对传统文化持辩证的态度。他既批判了传统文化中的糟粕，如封建礼教、迷信思想等，也肯定了其中的精华，如爱国精神、民族气节等。

在《拿来主义》一文中，鲁迅提出了"运用脑髓，放出眼光，自己来拿"的观点，主张对外来文化和传统文化都要采取批判继承的态度。这一观点对于我们今天理解和传承端午节等传统文化，仍然具有重要的指导意义。

\section{当代端午节的意义}

在当代社会，端午节的意义已经超越了传统的纪念和庆祝。它成为了：

\begin{itemize}
    \item \textbf{文化认同的象征}：端午节是中华民族文化认同的重要标志，增强了民族凝聚力。
    \item \textbf{家庭团聚的契机}：端午节为家人团聚提供了机会，增进了家庭感情。
    \item \textbf{文化传承的载体}：通过各种习俗活动，年轻一代了解和传承了传统文化。
    \item \textbf{爱国主义教育的平台}：通过纪念屈原，弘扬爱国主义精神。
\end{itemize}

\section{结语}

鲁迅先生虽然生活在动荡的年代，但他对传统文化的思考至今仍具有重要的现实意义。端午节作为中华民族的重要传统节日，承载着丰富的文化内涵和民族精神。

在当代，我们应该继承和发扬端午节的文化价值，摒弃其中的陋习，让这一传统节日焕发新的生机。同时，我们也应该学习鲁迅先生的批判精神，以辩证的态度对待传统文化，取其精华，去其糟粕，让传统文化在新时代焕发出更加绚丽的光彩。

端午节不仅是对历史的纪念，更是对未来的展望。让我们在庆祝端午节的同时，传承民族精神，弘扬中华文化，为实现中华民族的伟大复兴而努力。

\end{document}